\documentclass[10pt,a4paper]{amsart}
\usepackage[margin=19mm]{geometry}
\usepackage{amsmath,amssymb,mathtools}
\usepackage[scr=rsfs,cal=euler]{mathalfa}
\usepackage{xfrac}
\usepackage[unicode,hidelinks,bookmarks=false]{hyperref}
\newcommand{\ie}{i.e.\ }
\newcommand{\cf}{cf.\ }
\newcommand{\resp}{respectively\ }
\newcommand{\Subsection}{\subsection}
\providecommand{\nobreakdash}{\nobreak\hbox{-}}
\setlength{\emergencystretch}{2em}
\multlinegap0pt
\usepackage{amssymb}
\usepackage{tikz}
\newtheorem{theorem}{Theorem}
\newtheorem{lemma}{Lemma}
\title{The Graph Minor Relation Satisfies the Twin Alternative Conjecture}
\author{Jorge Bruno}
\date{Source: arXiv:2212.07670v2 (CC BY 4.0)}
\begin{document}
\maketitle

\noindent\emph{Source and licence.} The Graph Minor Relation Satisfies the Twin Alternative Conjecture. Version arXiv:2212.07670v2 (CC BY 4.0), distributed by arXiv under the Creative Commons Attribution 4.0 International licence, http://creativecommons.org/licenses/by/4.0/, as recorded in the arXiv licence metadata for this exact version and shown on the abstract page https://arxiv.org/abs/2212.07670v2. Full licence notice: \texttt{licenses/arxiv-2212.07670-CC-BY-4.0.txt}.\par\smallskip

\noindent\textit{Editorial scope.} Counted unit: the article's theorem theorem (source0.tex lines 253--254). The article's complete original proof of it is delivered with it (source0.tex lines 256--321, one \texttt{proof} environment of 66 lines). No mathematics is written, repaired or re-derived: the statement and the proof are the article's own lines, in the article's own notation, and no retained line is substituted or re-flowed.  No further source-local context is needed: every cross-reference of every retained line resolves inside the two delivered spans.  Nothing was omitted from any retained span, so the package carries no omission record.  No retained line contains a citation, so the wrapper carries no bibliography and no bibliography entry.  No retained line contains a figure environment, an image-inclusion command or drawing code, so no image is delivered and the package declares no asset dependency.  Editorial formatting only, in addition: an AMS \texttt{amsart} preamble that restates the article's own declarations and macros used by the retained lines, namely the environments \texttt{theorem}, \texttt{lemma} on separate counters, together with the standard packages those lines need.  The retained environments print in this document as Theorem 1, Lemma 1 in order of appearance, whereas the article prints the same items with the numbering of its own full document.  The complete native source of the article as distributed in the arXiv e-print for this exact version is bundled unchanged as \texttt{sources/134921/source0.tex}.
\begin{theorem} For any small tree $(T,r)$ with $|[(T_v, v)]_*| =1$, $\forall v \in succ(r)$, it follows that $|[(T,r)]_*| =1$ or $\infty$. 
\end{theorem}

\begin{proof} Let $(T,r)$ and $(S,s)$ be two non-isomorphic trees with $(S,s)\in [(T,r)]_*$, and models $\mu_T: V(T) \to \text{Sub}(S)$ and $\mu_S: V(S) \to \text{Sub}(T)$. Further assume that $|[(T_v, v)]_*| =1$ for any $v \in succ(r)$. We seek to prove that $|[(T,r)]_*| \geq \aleph_0$. 

Put $\mathcal{I}= \{(X_\alpha,\alpha) \mid \alpha\in \kappa\}$  as the set of all isomorphism types of all rooted trees $(T_v,v)$ and $(S_w,w)$ with $v\in succ(r)$  and $w\in succ(s)$. Define $f_T: \{(T_v,v)\mid v\in succ(r)\} \to \kappa$ so that $f_T(T_v,v) = \alpha$ provided $(T_v,v) \cong (X_\alpha, \alpha)$ (resp. $f_S: \{(S_w,w)\mid w \in succ(s)\} \to \kappa$ so that $f_S(S_w, w) = \alpha$ provided $(S_w, w) \cong (X_\alpha, \alpha)$). Since $(T,r) \not \cong (S,s)$ there must exist a $\beta \in \kappa$ for which, without loss of generality, $ |f_T^{-1}(\beta)| < |f_S^{-1}(\beta)|$. Put $T_B = f_T^{-1}(\beta)$ and $S_B = f_S^{-1}(\beta)$, and $\eta = |f_T^{-1}(\beta)|$ and $\zeta = |f_S^{-1}(\beta)|$. Enumerate the elements of $T_B$ and $S_B$ as $(T_i,i)$ and $(S_j,j)$ with $i\in \eta$ and $j\in \zeta$. We are then left with the following complementary scenarios.\\



\noindent
{\bf CASE 1:} $\forall j \in \zeta$, $\exists i_j \in \eta$ such that $V(\mu_S(S_j)) \subseteq V(T_{i_j})$\\ 

\noindent
{\bf CASE 2:} $\exists j \in \zeta$ with $V(\mu_S(S_\gamma)) \subseteq V(T_v)$, $(T_v, v) \not \in f_T^{-1}(\beta)$. \\

Before investigating the above cases we establish a simple yet useful result.

\begin{lemma}\label{lem:root_anchor} Let $(T_1, r_1)$ and $(T_2, r_2)$ be isomorphic small rooted trees. If $\mu: V(T_1) \to \text{Sub}(T_2)$ is a model witnessing $(T_1, r_1) \leq^* (T_2, r_2)$ such that $r_2 \notin V(\mu(r_1))$, then $(T_2, r_2)$ must be a bare ray or a finite path.\end{lemma}\begin{proof}Suppose $r_2 \notin V(\mu(r_1))$. Because $V(\mu(r_1))$ is a connected subset that does not contain the root $r_2$, the entire image of $T_1$ under $\mu$ must be mapped strictly "above" $r_2$, meaning it is contained within some proper principal subtree $(T_{2,v}, v)$ for a successor $v \in succ(r_2)$. Because $(T_1, r_1) \cong (T_2, r_2)$, this embedding implies that $(T_2, r_2)$ can be minor-embedded into a proper principal subtree of itself. Iterating this embedding generates an infinite sequence of vertices $r_2 <_T v = x_1 <_T x_2 <_T \ldots$, where each $x_n$ roots a nested isomorphic copy of the tree. This strictly increasing sequence of vertices spans a ray $R$ in $T_2$. Clearly, $R$ is a bare ray precisely when $(T_2, r_2)$ is a bare ray.\end{proof}

%Before moving further with this proof, let us establish some more notation: put $(T',r)$ as $(T,r)$ but with all full subtrees $(T_i,i)$, $i \in \eta$, removed - this also includes the roots $i$ of each $(T_i,i)$ and the edge joining $i$ with $r$. In what follows we fix a $k \in \eta$ and for each cardinal $\lambda$ denote $(T'_\lambda,r)$ as $(T',r)$ with $\lambda$ copies of $(T_k,k)$ attached to $r$ - we label each such copy $(T'_{\lambda_l},{\lambda_l})$, $l\in \lambda$. Clearly, $(T'_\eta,r) \cong (T,r)$ and $(T'_\nu,r) \not \cong (T'_\xi,r)$ for any pair of cardinals $\nu \not = \xi$. \\

\noindent
{\bf CASE 1:}
Observe that $\zeta = \sum_{i \in \eta} |A_i|$, where$$A_i = \{(S_j,j) \in f_S^{-1}(\beta) \mid V(\mu_S(S_j)) \subseteq V(T_i)\}.$$Because $\eta < \zeta$ there must exist at least one $\gamma \in \eta$ such that $|A_\gamma| \geq 2$. Recall that $(T_\gamma, \gamma) \cong (S_j, j)$ and for all $j \in \zeta$. The fact that $|A_\gamma| \geq 2$ implies that the subtree $(T_\gamma, \gamma)$ contains at least two disjoint models of itself. We can then apply Theorem~\ref{lem:root_anchor} to obtain a contradiction.\\


\noindent
{\bf CASE 2:} First, let $(T', r)$ be the rooted tree obtained from $(T, r)$ by deleting all $\eta$ principal subtrees that are isomorphic to $(S_\gamma, \gamma)$. That is, all trees isomorphic to $(S_\gamma, \gamma)$ rooted at a successor or $r$, along with their root and, thus, edge joining them to $r$. For any cardinal $\lambda$, we define the tree $(T'_\lambda, r)$ as the tree $(T', r)$ with exactly $\lambda$ new, disjoint copies of $(S_\gamma, \gamma)$ attached to the root $r$.\\

First consider the case where $\eta$ is finite (tacitly including the case where $\eta = 0$). Fix a $\gamma \in \zeta$ for which $\mu(S_\gamma,\gamma) \subseteq (T_v,v)\not \in f_T^{-1}(\beta)$. Next, for each $n\in \omega$, put $w_n \in succ(s)$ so that $(\mu_T\circ^*\mu_S)^n[(S_\gamma,\gamma)] \subseteq (S_{w_n}, w_n)$ and $v_n\in succ(r)$ with $\mu_S\circ^*(\mu_T\circ^*\mu_S)^n[(S_\gamma,\gamma)] \subseteq (T_{v_n}, v_n)$. Clearly, $(S_{w_0},w_0) = (S_\gamma,\gamma)$ and $(T_{v_0}, v_0) = (T_v, v)$, and 
\[
(S_{w_0},w_0) \leq^* (T_{v_0},v_0) \leq^*(S_{w_1},w_1) \leq^*(T_{v_1},v_1) \leq^* \ldots
\]
 Observe that $w_n \not = w_0$ for any $n\in \omega$ since then $(S_\gamma, \gamma) \cong (T_v,v)$ from the assumption that $|[(T_u, u)]_*| =1$ for any $u \in succ(r)$. In general, we claim that the sequences $(w_n)_{n \in \omega}$ and $(v_n)_{n \in \omega}$ consist of strictly non-repeating elements. Suppose for contradiction that the sequence loops. Without loss of generality, assume $w_1 = w_2$. The resulting mutual embedding forces $(S_{w_1}, w_1) \equiv^* (T_{v_1}, v_1)$. Under our assumption that $|[(T_u, u)]_*| = 1$ for all immediate successors, this equivalence guarantees strict isomorphism: $(S_{w_1}, w_1) \cong (T_{v_1}, v_1)$. Theorem~\ref{lem:root_anchor} states that $v_1 \in V(\mu_T(w_1))$. Since $\mu_T$ also maps $(T_{v_0},v_0)$ into $(S_{w_1}, w_1)$ and $v_0$, $v_1$ are incomparable in $T$, then $\mu_T$ would fail to preserve the meet semilattice order of $T$. Thus, we get two sequences of non-repeating elements: $(w_n)_{n\in\omega}$ and $(v_n)_{n\in\omega}$. \\

\noindent
We finish the case where $\eta$ is finite as follows. Clearly, $(T'_{\lambda}, r) \geq^* (T,r)$ for each $\lambda \geq \eta$. Working in the other direction, given any finite $\lambda \geq \eta$, let $u_1, \dots, u_\lambda \in succ(r)$ be the roots of the $\lambda$ new principal subtrees in $T'_\lambda$ that are isomorphic to $(S_\gamma, \gamma)$. We define the model $\mu_\lambda: V(T'_{\lambda}) \to \text{Sub}(T)$ piecewise as follows:
\begin{itemize}
\item For each new copy $1 \leq l \leq \lambda$, let 

\[
\mu_\lambda \restriction T'_{u_l} = \mu_S \circ^* (\mu_T \circ^* \mu_S)^{l-1},
\]

\item For the original chain branches $n \in \omega$, let 

\[
\mu_\lambda \restriction T_{v_n} = (\mu_S \circ^* \mu_T)^{\lambda},
\]

\item Let $\mu_\lambda$ act as the identity mapping everywhere else in $T'_\lambda$.
\end{itemize}

\noindent
Since $(T'_{u_l}, u_l) \cong (S_\gamma, \gamma)$, this mapping embeds the $\lambda$ new copies into the first $\lambda$ slots of the chain, while the operator $(\mu_S \circ^* \mu_T)^\lambda$ shifts the original branches downward without collision. Thus, $\mu_\lambda$ is a valid model.


The case for $\eta$ infinite follows a similar flavour to the finite case; we prove that $(T', r) \equiv^* (T,r)$ (and thus, $(T'_\lambda, r) \equiv^* (T,r)$ for all cardinals $\lambda \leq \eta$). Clearly, we need only show $(T, r) \leq^* (T',r)$. Since $\zeta > \eta$ (and both are infinite), then there exists a subset $S'_B \subset S_B$ with $|S'_B| = \zeta$ such that for every $(S_w, w) \in S'_B$, its image under $\mu_S$ is entirely contained within a strictly non-$\beta$-type principal subtree of $T$. Because $|T_B| = \eta < \zeta = |S'_B|$, there exists an injection $m: T_B \to S'_B$. Let $N_{recv}$ denote the set of non-$\beta$-type principal subtrees $(T_u, u)$ in $T$ that contain the image $\mu_S(m(T_v, v))$ for some $(T_v, v) \in T_B$. 

Finally, we define the model $\mu: V(T) \to \text{Sub}(T')$ piecewise as follows:
\begin{enumerate}
\item For each $(T_v, v) \in T_B$, let $\phi_v: T_v \to m(T_v)$ be an isomorphism. We set $\mu \restriction T_v = \mu_S \circ^* \phi_v$.
\item For each receiving non-$\beta$-type branch $(T_u, u) \in N_{recv}$, we evacuate the native infrastructure by setting $\mu \restriction T_u = (\mu_S \circ^* \mu_T)^2 \restriction T_u$.
\item Let $\mu$ act as the identity mapping on all remaining vertices of $(T, r)$.
\end{enumerate}

For the exact same reasons outlined in the finite case ($\eta \in \omega$), the shifted native infrastructure in Part (2) strictly does not clash with the incoming $\beta$-copies mapped in Part (1). This establishes that $\mu$ is a valid graph minor model, yielding $(T, r) \leq^* (T', r)$ and completing the proof.

\end{proof}

\end{document}
